%!TEX root = ../hutbthesis_main.tex
\chapter{系统总体设计}
\section{系统需求分析}
本课题面向方向盘视觉引导操控场景，目标不是单纯完成静态模型展示，而是形成能够加载人体上肢肌肉骨骼模型、驱动方向盘运动、提供实时反馈并完成精度评估的闭环仿真系统。因此，在总体设计阶段需要先明确系统的使用对象、任务边界、功能组成和性能约束，保证后续建模、算法实现和测试分析具有统一依据。
\begin{table}
  \centering
  \caption{系统需求分析}
  \label{tab:system_requirements}
  \begin{tabular}{ccc}
    \toprule
    需求类别 & 具体内容 & 设计对应关系 \\
    \midrule
     模型需求	& 加载上肢肌肉骨骼模型与方向盘模型，完成姿态匹配和物理绑定 & 对应仿真建模层与模型初始化流程 \\
     控制需求	& 将肌肉激活信号转换为平滑、受限、可执行的方向盘控制指令	& 对应控制算法层中的滤波、死区、增益和限幅处理 \\
     交互需求	& 实时显示目标角度、当前角度、误差变化和任务状态	& 对应视觉引导层与反馈界面 \\
     评估需求	& 采集实验数据并计算 MAE、RMSE、响应延迟、最大误差等指标	& 对应数据采集与评估层 \\
     工程需求	& 参数可配置、模块可替换、实验可复现，便于调试和论文验证	& 对应配置管理、模块化接口和结果保存机制 \\
    \bottomrule
  \end{tabular}
\end{table}


（1）模型需求系统需要载入上肢肌肉骨骼模型和方向盘模型，完成初始姿态的匹配和物理绑定，让两个模型能稳定结合，这部分需求对应仿真建模层和模型初始化的流程，为后续肌肉驱动和方向盘联动提供基础的物理环境支持。

（2）控制需求系统需要把肌肉激活信号转成平稳、受限、能被执行的方向盘控制指令，通过滤波、死区、增益和限幅等处理步骤，抑制噪声带来的干扰，避免出现超调和抖动问题，这部分需求对应控制算法层的信号处理模块，保证操控过程流畅，也能让系统保持可控状态。

（3）交互需求系统需要实时显示目标角度、当前角度、误差变化和任务状态，支持阶跃、正弦等目标轨迹的可视化展示，给操作人员提供清晰的视觉反馈，这部分需求对应视觉引导层和反馈界面，帮助用户直观了解操控状态和误差情况。

（4）评估需求系统需要收集仿真过程里的实验数据，包括时间戳、目标角度、实际角度、控制输出等内容，并且自动计算 MAE、RMSE、响应延迟、最大误差等指标，为性能分析提供可以参考的量化数据，这部分需求对应数据采集与评估层，为后续的精度验证和结果分析提供支撑。

（5）工程需求系统需要具备参数可配置、模块可替换、实验可复现的特点，支持模块化接口设计和结果保存的机制，方便调试、扩展和论文验证，这部分需求对应配置管理和整体架构设计，让系统的维护难度更低，也能更好地适配不同的使用场景。

综合以上这些需求，本系统主要需要解决三个核心问题，第一，如何让人体上肢模型和方向盘模型形成稳定的物理结合；第二，如何把肌肉激活信号转成连续、可控的方向盘转角指令；第三，如何通过视觉反馈和量化指标，证明系统的操控精度和运行稳定性达到了要求。


\section{系统架构设计}
根据系统的相关需求，本研究搭建了四层模块化的整体架构，这四层分别是仿真建模层、控制算法层、视觉引导层和数据采集与评估层，模型层主要提供物理运行环境，控制层负责生成对应的动作指令，视觉层用来给出交互反馈，评估层则负责输出系统的性能结果，四个层次按照这样的逻辑依次连接，这样的结构既能让系统的运行流程更加清晰，也方便在后续的工作中对单个模块进行调试和替换。

仿真建模层处在整个系统的最底部，主要用来完成人体上肢肌肉骨骼模型、方向盘模型和约束关系的搭建工作，这一层会通过 XML 模型文件对刚体结构、关节自由度、肌肉作动器、阻尼、摩擦和转角限制等物理属性进行说明，还会调用 MuJoCo 动力学求解器完成运动状态的更新，这一层输出的人体关节状态、肌肉激活水平和方向盘转角，都是控制算法层和评估层需要用到的基础数据。

控制算法层处在系统的核心位置，主要负责将肌肉激活信号或者上层控制信号转换成方向盘可以执行的指令，这一层包含信号预处理、低通滤波、死区过滤、非线性映射、增益调节和转角限幅等处理步骤，通过这些处理步骤，系统可以降低小幅噪声带来的方向盘抖动问题，避免大幅信号带来的超调情况，同时也能让方向盘在目标角度附近具备更好的细微调节能力。

视觉引导层主要服务于操作人员和实验观察人员，它会生成目标轨迹并实时展示系统当前的运行状态，这一层会把阶跃轨迹、正弦轨迹等目标角度叠加显示在仿真界面上，同时也会呈现方向盘的当前角度、瞬时误差和任务完成状态，操作人员可以直观看到目标、实际角度和误差的变化情况；数据采集与评估层会在仿真运行的过程中同步记录时间戳、目标角度、实际角度、控制输出和误差序列，还会在测试结束后计算控制精度、响应延迟和运行稳定性等相关指标，图~\ref{fig:system_architecture}展示系统这四层架构之间的联系。

从数据传输的方向来看，系统会先由仿真建模层提供可以观测的状态，控制算法层根据当前误差和肌肉激活情况生成对应的控制指令，视觉引导层将目标和反馈结果展示给用户，评估层再将运行过程转换成可以量化的实验数据，这样的设计方式可以避免单个模块承担太多的功能，也可以让第四章的建模实现、第五章的算法实现和第六章的测试分析形成前后对应的关系。

\begin{figure}[htbp]
  \centering
  \includegraphics[width=0.8\textwidth]{image_5.png}
  \caption{系统整体架构}
  \label{fig:system_architecture}
\end{figure}






\section{系统工作流程}
系统的工作流程可以分成初始化、目标生成、闭环控制、数据记录和结果输出这五个阶段，初始化阶段主要完成配置读取、模型载入、关节限位检查、方向盘初始角度归零和仿真窗口启动，让人体模型、方向盘模型和控制参数都保持统一的状态，这个阶段完成之后，系统就会进入正式的仿真循环。目标生成阶段会按照测试任务生成需要的方向盘角度，要是开展定点转向测试，系统就会生成阶跃目标角度，用来查看方向盘能不能快速到达并稳定保持在目标位置，要是开展连续跟踪测试，系统就会生成正弦目标轨迹，用来检验系统对动态目标的跟踪效果，目标角度确定以后，系统会同步读取当前方向盘角度和肌肉激活状态，并且计算出实时的角度误差。

闭环控制阶段是系统运行的重要部分，控制算法先对输入信号做滤波和死区处理，再根据误差大小完成增益调整和非线性映射，最后输出方向盘的控制指令，MuJoCo 接收到控制指令后会推进一个仿真步长，更新人体关节、肌肉和方向盘的运行状态，视觉引导界面也会同步刷新目标角度、实际角度和误差提示，形成感知、决策、控制、反馈的完整闭环，图~\ref{fig:system_workflow}展示了系统的整体工作流程。

数据记录与结果输出会贯穿整个测试过程，系统在每一个采样时刻都会保存目标角度、实际角度、角度误差、控制输出和仿真时间，当达到设定的测试时长或者用户主动停止时，系统就会结束仿真并输出评估结果，这些结果包括 MAE、RMSE、响应延迟、最大误差、漂移量和任务完成率等指标，通过这样的流程，系统不仅可以完成方向盘操控的演示，还能为论文中的精度分析提供可以重复验证的数据依据。

\begin{figure}[htbp]
  \centering
  \includegraphics[width=0.8\textwidth]{image_6.png}
  \caption{系统工作流程图}
  \label{fig:system_workflow}
\end{figure}  
















\section{开发环境与工具}
本系统在 Windows 10 64 位操作系统下完成开发与测试，仿真核心采用 MuJoCo 2.3.3 多刚体动力学引擎，项目框架基于 User-in-the-Box 开源框架搭建。开发语言为 Python 3.8，主要使用 NumPy 完成数值计算与数组处理，使用 Matplotlib 完成实验数据绘图和结果可视化，使用 mujoco Python 接口完成模型加载、仿真推进、状态读取与控制指令下发。

系统使用的开发工具为 VS Code，主要用来完成代码编写、模块调试和工程文件管理工作，模型文件使用 XML 格式进行描述，配置文件使用 YAML 格式保存模型参数、控制参数和测试参数，在工程结构上，系统按照功能分成模型资源、控制算法、视觉反馈、数据记录和结果评估等不同模块，模型资源模块用于存放人体上肢模型、方向盘模型以及相关的材质文件，控制算法模块负责处理肌肉信号和完成转角映射，视觉反馈模块负责显示目标轨迹和误差情况，数据记录模块用于保存实验过程中的相关数据，结果评估模块负责计算控制精度和稳定性指标，不同模块通过统一的接口进行数据交换，避免直接依赖具体的文件路径，提升了系统复现和维护的方便程度。通过以上开发环境和工程组织方式，本系统可以在普通的 PC 设备上完成肌肉驱动方向盘操控的仿真实验，相关的模型、参数和测试结果都可以在项目工程内部进行统一管理，既可以满足毕业设计实验验证的需求，也方便后续在不同设备上进行重新部署和功能扩展。
